\subsection{Shared affine preliminaries}

We first record the support, relabelling, and replacement facts that let us
move between comparison environments.  They require no continuity or
compounding axiom.

Recall that $(q,K)\succeq(p,L)$ denotes the paper's canonical two-block
comparison.  For channels $E,G$ at a common prior $q$, write
$E\succeq_q G$ for $(q,E)\succeq(q,G)$.  A \emph{finite experiment at $q$}
is a finite payoff--record channel paired with that prior; its marked terminal
law is denoted by $\eta_E$.

\begin{lemma}[Blocks, supports, and dummies]\label{lf:lem:bookkeeping}
Suppose Axioms \textup{(A1)} and \textup{(A5)--(A7)} hold.
\begin{enumerate}[label=\textup{(\alph*)},leftmargin=*]
\item After identifying alternatives related by invertible action or record
  relabellings, the derived comparison of pairs is complete and transitive.
  If either member of a pair comparison is replaced by an indifferent
  alternative, the comparison is unchanged.
\item For every finite experiment $(q,K)$, if $S=\supp q$, then $(q,K)$ is indifferent to the restriction of
  $(q,K)$ to $S$.
\item For every finite experiment $(q,K)$, every nonempty finite $B$, and
  every $s\in\D(B)$, if $K^\uparrow(o,r\mid a,b):=K(o,r\mid a)$, then
  \begin{equation}
  (q\otimes s,K^\uparrow)\sim(q,K).
  \label{lf:eq:dummy-ordinal}
  \end{equation}
\end{enumerate}
\end{lemma}

\begin{proof}
Place any finite collection of alternatives in one finite block environment.
Axiom~\textup{(A5)} removes unused blocks, and Axiom~\textup{(A1)} supplies
completeness and transitivity there.  Invertible action and record relabellings
are neutral by applying Axioms~\textup{(A6)} and~\textup{(A7)} in both
directions.  In particular, swapping both tags identifies the two orientations
of a two-block comparison.  This proves part \textup{(a)}; a three-block
environment gives transitivity and a four-block environment gives replacement.

For part \textup{(b)}, process an action in $S$ by its inclusion in $A$.  In the
other direction, output $a\in S$ when the realised action is $a$, and choose an
arbitrary element of $S$ off the support.  Both processors satisfy
\eqref{eq:action-processing}; rows attached to null processed actions may be
completed arbitrarily.  Axiom~\textup{(A7)} gives the two comparisons.

For part \textup{(c)}, first project $(a,b)$ to $a$.  Conversely, given $a$,
draw $b$ independently from $s$ and output $(a,b)$.  Again the two joint-law
identities in \eqref{eq:action-processing} are exact, so Axiom~\textup{(A7)}
applies in both directions.  Part \textup{(a)} completes both arguments.
\end{proof}

Fix for the moment a full-support prior $q\in\D(A)$.  For a joint channel $K$,
write
\[
p_K(o,r):=\sum_a q(a)K(o,r\mid a),
\qquad
\pi^K_{or}(a):=\frac{q(a)K(o,r\mid a)}{p_K(o,r)}
\]
when $p_K(o,r)>0$, and define its \emph{marked terminal law}
\begin{equation}
\eta_K:=\sum_{(o,r):p_K(o,r)>0}
p_K(o,r)\,\delta_{(o,\pi^K_{or})}.
\label{lf:eq:marked-law}
\end{equation}
It records the terminal payoff and the posterior about the realised action.
Write $I_q(E)$ for the mutual information between the action and the complete
visible output of $E$ under $q$.  When the output is a first-stage record $Y$,
we also write this quantity as $I_q(A;Y)$.

\begin{lemma}[Terminal laws and affine fibres]\label{lf:lem:terminal-affine}
Under Axioms \textup{(A1)}, \textup{(A2)}, and
\textup{(A5)--(A8)}, the following statements hold.
\begin{enumerate}[label=\textup{(\alph*)},leftmargin=*]
\item If $q$ has full support and $K,L$ are finite experiments at $q$ with
  $\eta_K=\eta_L$, then
  $(q,K)\sim(q,L)$.
\item At every full-support $q$, the quotient of finite experiments by their
  marked terminal laws has an affine representation $W_q$.
\end{enumerate}
\end{lemma}

\begin{proof}
We give the finite Blackwell argument behind part \textup{(a)}.  For a payoff
$o$ and posterior $\pi$, let
\[
M_K(o,\pi):=
\sum_{\substack{r:\,p_K(o,r)>0\\ \pi^K_{or}=\pi}} p_K(o,r).
\]
Equality of marked laws gives $M_K(o,\pi)=M_L(o,\pi)=:M(o,\pi)$.
For $p_K(o,r)>0$, draw a record $s$ of $L$ according to
\begin{equation}
T(s\mid o,r):=
\1_{\{p_L(o,s)>0,\,\pi^L_{os}=\pi^K_{or}\}}
\frac{p_L(o,s)}{M(o,\pi^K_{or})}.
\label{lf:eq:matching-processor}
\end{equation}
Choose an arbitrary probability row when $p_K(o,r)=0$.  This is a
payoff-preserving record processor.  Bayes' formula gives
\[
q(a)K(o,r\mid a)=p_K(o,r)\pi^K_{or}(a).
\]
For a reached target record $s$, with posterior $\pi=\pi^L_{os}$, it follows
that
\[
\sum_r K(o,r\mid a)T(s\mid o,r)
=\frac{p_L(o,s)}{M(o,\pi)}
  \frac{\pi(a)}{q(a)}M(o,\pi)
=L(o,s\mid a).
\]
If $s$ is unreached, full support makes both sides zero.  Thus postprocessing
$K$ by $T$ gives $L$.  The symmetric construction gives a
processor from $L$ to $K$.  Axiom~\textup{(A6)} and
Lemma~\ref{lf:lem:bookkeeping}\textup{(a)} give indifference.

Boundary priors are handled later by first restricting to their positive
support using Lemma~\ref{lf:lem:bookkeeping}\textup{(b)}.

For part \textup{(b)}, take the quotient of actual finite marked experiments by
equality of \eqref{lf:eq:marked-law}.  If $E$ and $G$ are two such experiments,
their public mixture first draws a coin independent of the action, runs the
selected experiment, and retains the coin in the record.  Its marked law is
\begin{equation}
\eta_{tE\oplus(1-t)G}=t\eta_E+(1-t)\eta_G.
\label{lf:eq:public-mixture}
\end{equation}
Axiom~\textup{(A8)} gives, for $0<t<1$ and any third experiment $L$,
\begin{equation}
E\succeq_q L
\quad\Longleftrightarrow\quad
tE\oplus(1-t)G\succeq_q tL\oplus(1-t)G.
\label{lf:eq:mixture-independence}
\end{equation}
This is the action-independent special case of the axiom: the public coin is the binary first-stage record and
$G$ is the common continuation.  If record alphabets differ, put them in a common disjoint union before invoking
the axiom; part \textup{(a)} identifies the result with
\eqref{lf:eq:public-mixture}.

It remains only to check the Archimedean step.  A closed segment between two
experiments, and any fixed comparator, can be realised on one fixed union of
their finite record alphabets.  Its channel entries vary continuously with
$t\in[0,1]$.
Axiom~\textup{(A2)} therefore makes both contour sets along the segment closed.
Completeness says that they cover $[0,1]$; when a target lies between the two
endpoints, the sets contain opposite endpoints and connectedness makes them
intersect.  Thus every such target is indifferent to a point of the segment.
The mixture-space theorem now supplies an affine representative $W_q$.  This
quotient is closed under every mixture and embedding used below.
\end{proof}

\subsection{From the axioms to their working forms}

The axioms were stated in their weakest forms: a binary sure-thing principle and fixed-channel relevance
comparisons.  The proof uses their working forms: finite-branch continuation and environment-comparison
relevance.  This subsection records the two equivalences and makes explicit where the other axioms enter
them.  Only the statements are used below.

\begin{lemma}[Finite-branch extension]\label{pt:lem:finite-branch-extension}
Suppose Axioms~\textup{(A1)} and \textup{(A5)--(A7)} hold.  Then
Axiom~\textup{(A8)} is equivalent to the following property.  Given
$k\ge1$, a first-stage record channel
$P:A\to\D(\{1,\ldots,k\})$, a prior $q\in\D(A)$, and two ordered lists of
continuation channels $(K_1,\ldots,K_k)$ and $(L_1,\ldots,L_k)$, with
$K_i,L_i:A\to\D(O\times R)$, define for every reached $i$
\[
 m_i:=\sum_aq(a)P(i\mid a),
 \qquad q_i(a):=\frac{q(a)P(i\mid a)}{m_i}.
\]
If every reached $i$ satisfies
\[
 (q_i,K_i)\succeq(q_i,L_i),
\]
then
\begin{equation}
 \bigl(q,P\triangleright(K_1,\ldots,K_k)\bigr)
 \succeq
 \bigl(q,P\triangleright(L_1,\ldots,L_k)\bigr),
 \label{pt:eq:finite-branch-extension}
\end{equation}
strictly if one reached branch comparison is strict.
\end{lemma}

\begin{proof}
Assume first Axiom~\textup{(A8)}.  Put
\[
 m_i:=\sum_aq(a)P(i\mid a),
 \qquad
 I^+:=\{i\in\{1,\ldots,k\}:m_i>0\}.
\]
If $m_i=0$, then $P(i\mid a)=0$ for every $a\in\supp q$.  Consequently, changing the continuation after such a
record changes the compound channel only on action rows outside $\supp q$.  Apply Axiom~\textup{(A7)} in both
directions using the identity action processor: its Bayesian pushforwards agree on supported rows and allow
arbitrary completions elsewhere.  Continuations following unreached records are therefore immaterial.  We may
hold them fixed and enumerate the reached records as $I^+=\{i_1,\ldots,i_n\}$.

If $n=1$, the first-stage record is constant on $\supp q$.  Deleting that
constant label and adjoining it again are payoff-preserving record processings
in both directions on every reached row; the off-support rows are immaterial
by the preceding paragraph.  Axiom~\textup{(A6)} therefore identifies the
compound comparison with the sole branch comparison.  Suppose henceforth that
$n\ge2$.

For $j=0,\ldots,n$, define hybrid continuation channels by
\[
 H_i^j:=
 \begin{cases}
 K_i,&i\in\{i_1,\ldots,i_j\},\\
 L_i,&i\in I^+\setminus\{i_1,\ldots,i_j\},
 \end{cases}
\]
and use a common arbitrary continuation on $\{1,\ldots,k\}\setminus I^+$.  Write
$E_j:=P\triangleright(H_1^j,\ldots,H_k^j)$.  Fix $j\ge1$ and collapse the first-stage record to the binary channel
\[
 B_j(1\mid a):=P(i_j\mid a),
 \qquad
 B_j(2\mid a):=\sum_{i\ne i_j}P(i\mid a).
\]
Both binary records are reached.  Put all binary continuations on the common record set
$\{1,\ldots,k\}\times R$.  Pad the
two continuations that may differ by
\[
 \widetilde K_j(o,(i,r)\mid a):=\1_{\{i=i_j\}}K_{i_j}(o,r\mid a),
 \qquad
 \widetilde L_j(o,(i,r)\mid a):=\1_{\{i=i_j\}}L_{i_j}(o,r\mid a).
\]
On record $2$, collect all unchanged branches in the continuation channel
\begin{equation}
 M_j(o,(i,r)\mid a)
 :=\frac{P(i\mid a)}{B_j(2\mid a)}H_i^j(o,r\mid a),
 \qquad i\ne i_j,
 \label{pt:eq:complement-continuation}
\end{equation}
with zero probability on records $(i_j,r)$, whenever $B_j(2\mid a)>0$, and complete the other rows arbitrarily.
Since $H_i^j=H_i^{j-1}$ for $i\ne i_j$, this continuation is common to the two comparisons.
Binary record $1$ has probability
$m_{i_j}$ and carries the conditional action lottery $q_{i_j}$; record $2$ carries the complementary conditional
lottery, and their mixture is $q$.  On the active record support, deleting or
adjoining the constant label $i_j$ gives payoff-preserving record processings
in both directions.  Axiom~\textup{(A6)} therefore identifies the comparison
of $\widetilde K_j$ and $\widetilde L_j$ with that of $K_{i_j}$ and $L_{i_j}$.
Hence Axiom~\textup{(A8)} gives
\[
 \bigl(q,B_j\triangleright(\widetilde K_j,M_j)\bigr)
 \succeq
 \bigl(q,B_j\triangleright(\widetilde L_j,M_j)\bigr),
\]
strictly whenever the comparison at $q_{i_j}$ is strict.

On the active record support, the nested records are in bijection with the
flat records through
\[
 (1,(i_j,r))\longmapsto(i_j,r),
 \qquad
 (2,(i,r))\longmapsto(i,r)\quad(i\ne i_j).
\]
Extend these maps arbitrarily over records of zero probability.  Together
with the preceding removal of off-support rows, they give payoff-preserving
record processings in both directions.  Axiom~\textup{(A6)} therefore
makes the two binary compounds indifferent to $E_j$ and $E_{j-1}$, respectively.  Thus
$(q,E_j)\succeq(q,E_{j-1})$, strictly at any strict replacement.  Place the finite collection
$E_0,\ldots,E_n$ in one block environment.  Axiom~\textup{(A5)} transports each pairwise comparison to that
environment, and transitivity in Axiom~\textup{(A1)} yields $(q,E_n)\succeq(q,E_0)$, strictly if at least one
step is strict.  Restoring the immaterial null branches proves \eqref{pt:eq:finite-branch-extension}.

Conversely, assume the displayed finite-branch property, including its strict clause.  In the binary environment
of Axiom~\textup{(A8)}, reflexivity and duplication coherence make the common continuation weakly comparable to
itself, so the finite-branch property gives the forward implication in
\eqref{eq:recordwise-sure-thing}.  If the aggregate comparison in that equation held while
$(q_1,K)\not\succeq(q_1,L)$, completeness would give $(q_1,L)\succ(q_1,K)$.  Applying the strict
finite-branch property to the reversed lists would give the strict reverse aggregate comparison, a contradiction.  This
proves the converse implication and hence Axiom~\textup{(A8)}.
\end{proof}

\begin{corollary}[One-branch insertion]\label{pt:cor:one-branch-a8}
Suppose Axioms~\textup{(A1)} and \textup{(A5)--(A8)} hold.  Fix $q\in\D(A)$, a first-stage record channel
$P:A\to\D(\{1,\ldots,k\})$, a reached record $i^*$ with posterior $q_{i^*}$, and continuation channels
$K_1,\ldots,K_k,L_1,\ldots,L_k:A\to\D(O\times R)$ such that $K_i=L_i$ for every $i\ne i^*$.  Then
\[
 (q_{i^*},K_{i^*})\succeq(q_{i^*},L_{i^*})
 \quad\Longleftrightarrow\quad
 \bigl(q,P\triangleright(K_1,\ldots,K_k)\bigr)
 \succeq
 \bigl(q,P\triangleright(L_1,\ldots,L_k)\bigr),
\]
and the same equivalence holds for strict preference.
\end{corollary}

\begin{proof}
The forward implication follows from Lemma~\ref{pt:lem:finite-branch-extension}, using duplication coherence on
the common branches.  If the aggregate comparison held but the branch comparison failed, completeness would make
the reversed branch comparison strict; the strict part of the finite-branch extension would then give a strict
reverse aggregate comparison.  The strict equivalence follows by applying the weak equivalence in both directions.
\end{proof}

In the appendices, references to Axiom~\textup{(A8)} with more than two first-stage records mean its
finite-branch extension in Lemma~\ref{pt:lem:finite-branch-extension}; when only one branch changes, we cite
Corollary~\ref{pt:cor:one-branch-a8}.

The relevance axioms in the main text are stated inside fixed channels.  The
proofs below use their equivalent comparison-environment forms.

For $o\in O$, write $K_o:\{\ast\}\to\D(O\times\{\ast\})$ for the channel that
delivers $o$ for sure.  On a finite action set $A$, write
$K^{\mathrm{id}}_o:A\to\D(O\times A)$ for the channel that delivers $o$ and
reports the action, and $K^0_o:A\to\D(O\times\{\ast\})$ for the channel that
delivers $o$ and no record.

\begin{lemma}[Fixed-channel and environment relevance]\label{lem:relevance-bridge}
Given Axioms~\textup{(A1)}, \textup{(A5)}, \textup{(A6)},
\textup{(A7)}, and \textup{(A8)}, the following equivalences hold.
\begin{enumerate}[label=\textup{(\alph*)},leftmargin=*]
\item Axiom~\textup{(A3)} is equivalent to the existence of
  $o^+,o^-\in O$ such that
  \[
    (\delta_\ast,K_{o^+})\succ(\delta_\ast,K_{o^-}).
  \]
\item Axiom~\textup{(A4)} is equivalent to the following condition: there is
  an outcome $o_\ast\in O$ such that, for every finite $A$ with $|A|\ge2$
  and every full-support $q\in\D(A)$,
  \[
    (q,K^{\mathrm{id}}_{o_\ast})\succ(q,K^0_{o_\ast}).
  \]
\end{enumerate}
\end{lemma}

\begin{proof}
For part \textup{(a)}, action processing in both directions identifies each
pure lottery in Axiom~\textup{(A3)} with the corresponding singleton
experiment; the unrestricted row at the unreached action supplies the reverse
pushforward.  Replacement transports the strict comparison both ways.

For part \textup{(b)}, restrict the witnessing lotteries to their supports.
Two-way action processing identifies them with fair binary full revelation
and a channel with two identical rows; two-way record processing identifies
the latter with silence.  Thus Axiom~\textup{(A4)} is equivalent to
\begin{equation}
  ((\tfrac12,\tfrac12),K^{\mathrm{id}}_{o_\ast})
  \succ
  ((\tfrac12,\tfrac12),K^0_{o_\ast})
  \label{eq:bridge-fair-binary}
\end{equation}
on a binary action set.

For any full-support binary $p$, choose
$0<\varepsilon\le2\min_i p_i$ and a record $y$ with
$P(y\mid i)=\varepsilon/(2p_i)$.  It has mass $\varepsilon$ and fair
posterior.  Reveal only after $y$ in one continuation plan and remain silent
everywhere in the other.  Equation~\eqref{eq:bridge-fair-binary} and
Corollary~\ref{pt:cor:one-branch-a8} rank the compounds strictly.  Full
revelation garbles to the first and the second garbles to silence, so
$(p,K^{\mathrm{id}}_{o_\ast})\succ(p,K^0_{o_\ast})$.

For general full-support $q$ on $|A|\ge2$, report a two-cell partition.
Two-way action processing identifies this with binary full revelation and
identifies silence with binary silence.  Hence the cell report is strictly
better than silence; full revelation weakly dominates the cell report by
record processing.  The converse takes the fair binary prior and reverses
the same support identifications.
\end{proof}
